\section{Introducción}
\label{sec:introduccion}

La anemia infantil representa un problema de salud pública crítico en las zonas rurales del Perú. Según la Encuesta Demográfica y de Salud Familiar (ENDES)~\cite{inei2023endes}, la prevalencia de anemia en niños de 6 a 35 meses alcanza el 44.7\% en áreas rurales bajo las últimas directrices, superando significativamente el promedio nacional. Los factores determinantes incluyen la pobreza, el limitado acceso a servicios de salud, la deficiente calidad del agua y la falta de educación nutricional en las comunidades andinas y amazónicas.

Los métodos tradicionales de diagnóstico de anemia requieren análisis de sangre en laboratorios clínicos, lo que implica desplazamientos largos y costosos para las familias rurales. Además, el seguimiento nutricional y la prevención de recaídas se ven obstaculizados por la falta de herramientas digitales que operen en entornos sin conectividad a internet.

En los últimos años, los sistemas de inteligencia artificial han demostrado un gran potencial en el diagnóstico no invasivo de enfermedades a través del análisis de imágenes. Particularmente, las redes neuronales convolucionales aplicadas a la evaluación de la conjuntiva ocular han mostrado resultados prometedores para la estimación de niveles de hemoglobina, permitiendo el diagnóstico en tiempo real mediante cámaras de teléfonos móviles~\cite{zhao2024prediction, uvadar2023NoninvasiveHE}. Sin embargo, la mayoría de estas soluciones requieren conexión a internet y no están adaptadas al contexto rural peruano.

Este trabajo presenta un sistema móvil con inteligencia artificial offline para el diagnóstico y prevención de la anemia infantil en zonas rurales del Perú. El sistema integra tres módulos principales: diagnóstico por análisis de conjuntiva ocular, recomendación nutricional personalizada basada en alimentos locales, y evaluación de riesgo familiar. La aplicación funciona completamente sin conexión a internet y utiliza una base de datos local de insumos alimenticios construida a partir de las Tablas Peruanas de Composición de Alimentos (TPCA)~\cite{INS2023} y mapeada a las 11 ecorregiones peruanas de Antonio Brack Egg~\cite{brack1986ecorregiones}. Los datos pueden sincronizarse mediante redes LoRa Mesh cuando estén disponibles.

El artículo está organizado de la siguiente manera. La Sección~\ref{sec:trabajos_relacionados} presenta los trabajos relacionados. La Sección~\ref{sec:materiales_metodos} describe la arquitectura del sistema y cada uno de los módulos. La Sección~\ref{sec:resultados} muestra los resultados obtenidos. Finalmente, las Secciones~\ref{sec:discusion} y~\ref{sec:conclusiones} presentan la discusión y las conclusiones del trabajo.

